\chapter{Symmetry Properties of the Fourier Transform}
\label{ch:symmetry}

In this chapter we study the interaction between the Fourier transform and the geometric symmetries of Euclidean space $E_n$, primarily rotations and dilations. A linear transformation on $E_n$ preserving the inner product is an orthogonal transformation (rotation or reflection). The group of rotations acts naturally on functions by $f \mapsto f \circ R$.

Since the Fourier transform commutes with rotations, $\widehat{f \circ R} = \hat{f} \circ R$, it is natural to decompose function spaces such as $L^2(E_n)$ into subspaces invariant under rotations and to examine how the Fourier transform acts on each invariant component. In $E_2$, this corresponds to the classical Fourier series expansion in polar coordinates $f(r e^{i\theta}) = \sum_{k=-\infty}^\infty f_k(r) e^{ik\theta}$. In $n$ dimensions, the role of $e^{ik\theta}$ is played by \emph{spherical harmonics}.

\section{The Decomposition of \texorpdfstring{$L^2(E_2)$}{L2(E2)} into Subspaces Invariant Under the Fourier Transform}

Let $x = (x_1, x_2) = (r \cos\theta, r \sin\theta) \in E_2$. Every function $f \in L^2(E_2)$ can be expanded in polar coordinates as $f(x) = \sum_{k=-\infty}^\infty f_k(r) e^{ik\theta}$, where $f_k(r) = \frac{1}{2\pi} \int_0^{2\pi} f(r\cos\theta, r\sin\theta) e^{-ik\theta}\,d\theta$, with convergence in $L^2(E_2)$. Let $\mathfrak{H}_k$ denote the subspace of $L^2(E_2)$ consisting of functions of the form $f(x) = f_k(r) e^{ik\theta}$; equivalently, $\mathfrak{H}_k$ consists of the functions $f$ such that $f(re^{i\theta}) e^{-ik\theta}$ is independent of $\theta$.

\begin{theorem}\label{thm:c4-1-1}
The subspaces $\mathfrak{H}_k$, $k \in \mathbb{Z}$, are mutually orthogonal and are invariant under rotations and under the Fourier transform.
\end{theorem}

\begin{corollary}\label{cor:c4-1-2}
If $f \in \mathfrak{H}_k$, then its Fourier transform $\hat{f}$ also belongs to $\mathfrak{H}_k$.
\end{corollary}

The following three relations summarize the basic facts about the decomposition and the action of the Fourier transform. First of all,
\begin{equation}\label{eq:c4-1-3}
L^2(E_2) = \bigoplus_{k=-\infty}^\infty \mathfrak{H}_k,
\end{equation}
the sum being orthogonal, and
\begin{equation}\label{eq:c4-1-4}
\|f\|_2^2 = 2\pi \sum_{k=-\infty}^\infty \int_0^\infty |f_k(r)|^2 r\,dr.
\end{equation}
Finally, if $f(x) = f_k(r) e^{ik\theta}$, then
\begin{equation}\label{eq:c4-1-5}
\hat{f}(\rho, \varphi) = (-i)^k 2\pi e^{ik\varphi} \int_0^\infty f_k(r) J_k(2\pi r \rho) r\,dr,
\end{equation}
where $J_k$ is the Bessel function of order $k$.

\begin{theorem}\label{thm:c4-1-6}
The Hankel transform of order $k$, defined by
\[
(H_k g)(\rho) = 2\pi \int_0^\infty g(r) J_k(2\pi r \rho) r\,dr,
\]
is a unitary operator on $L^2((0, \infty), r\,dr)$ and satisfies $H_k^{-1} = H_k$.
\end{theorem}

\section{Spherical Harmonics}

Let $\mathcal{P}_k$ denote the vector space of homogeneous polynomials of degree $k$ in $n$ real variables $x = (x_1, \dots, x_n)$. A polynomial $P \in \mathcal{P}_k$ is \emph{harmonic} if $\Delta P = \sum_{j=1}^n \partial^2 P / \partial x_j^2 = 0$. The space of all homogeneous harmonic polynomials of degree $k$ on $E_n$ is denoted by $\mathcal{H}_k$.

The restriction of a homogeneous harmonic polynomial $P \in \mathcal{H}_k$ to the unit sphere $\Sigma = \{x \in E_n : |x| = 1\}$ is called a \emph{spherical harmonic of degree $k$}. The space of all spherical harmonics of degree $k$ is denoted by $\mathcal{Y}_k$.

\begin{theorem}\label{thm:c4-2-1}
The dimension of the space $\mathcal{H}_k$ (and therefore of $\mathcal{Y}_k$) is given by
\[
d_k = \dim \mathcal{H}_k = \binom{n + k - 1}{k} - \binom{n + k - 3}{k - 2}.
\]
\end{theorem}

\begin{corollary}\label{cor:c4-2-2}
For $n = 2$, $d_0 = 1$ and $d_k = 2$ for all $k \geq 1$. For $n = 3$, $d_k = 2k + 1$.
\end{corollary}

\begin{corollary}\label{cor:c4-2-3}
Every homogeneous polynomial $P \in \mathcal{P}_k$ can be uniquely decomposed in the form $P(x) = P_k(x) + |x|^2 P_{k-2}(x) + |x|^4 P_{k-4}(x) + \dots$, where $P_{k-2j} \in \mathcal{H}_{k-2j}$.
\end{corollary}

\begin{corollary}\label{cor:c4-2-4}
The restriction of polynomials to the sphere $\Sigma$ is dense in $L^2(\Sigma)$; consequently $L^2(\Sigma) = \bigoplus_{k=0}^\infty \mathcal{Y}_k$.
\end{corollary}

Explicitly, the decomposition of Corollary~\ref{cor:c4-2-3} reads
\begin{equation}\label{eq:c4-2-5}
P(x) = P_k(x) + |x|^2 P_{k-2}(x) + |x|^4 P_{k-4}(x) + \dots,
\end{equation}
and the orthogonal decomposition of Corollary~\ref{cor:c4-2-4} is
\begin{equation}\label{eq:c4-2-6}
L^2(\Sigma) = \bigoplus_{k=0}^\infty \mathcal{Y}_k.
\end{equation}
Moreover, the spaces $\mathcal{Y}_k$ are mutually orthogonal:
\begin{equation}\label{eq:c4-2-7}
\int_\Sigma Y_k(t') \overline{Y_j(t')}\,dt' = 0 \quad \text{whenever } k \neq j.
\end{equation}

\begin{lemma}\label{lem:c4-2-8}
If $Y_k \in \mathcal{Y}_k$ and $Y_j \in \mathcal{Y}_j$ with $k \neq j$, then $Y_k$ and $Y_j$ are orthogonal in $L^2(\Sigma)$.
\end{lemma}

\begin{corollary}\label{cor:c4-2-9}
For each $k \geq 0$, let $\{Y_k^{(1)}, Y_k^{(2)}, \dots, Y_k^{(d_k)}\}$ be an orthonormal basis of $\mathcal{Y}_k$. The zonal harmonic of degree $k$ with pole at $\eta \in \Sigma$ is defined by
\[
Z_\eta^{(k)}(\xi) = \sum_{j=1}^{d_k} Y_k^{(j)}(\xi) \overline{Y_k^{(j)}(\eta)}.
\]
Then $Z_\eta^{(k)}$ depends only on the inner product $\xi \cdot \eta$.
\end{corollary}

\begin{theorem}\label{thm:c4-2-10}
Let $\lambda = (n-2)/2$. There exists a family of orthogonal polynomials $C_k^\lambda(t)$ on $[-1, 1]$ (the Gegenbauer or ultraspherical polynomials) such that
\[
Z_\eta^{(k)}(\xi) = \frac{d_k}{\omega_{n-1}} \frac{C_k^\lambda(\xi \cdot \eta)}{C_k^\lambda(1)}.
\]
\end{theorem}

\begin{lemma}\label{lem:c4-2-11}
The generating function for the Gegenbauer polynomials is
\[
(1 - 2rt + r^2)^{-\lambda} = \sum_{k=0}^\infty C_k^\lambda(t) r^k, \quad |r| < 1, |t| \leq 1.
\]
\end{lemma}

\begin{theorem}\label{thm:c4-2-12}
The Poisson kernel for the unit ball in $E_n$ can be expanded in terms of spherical harmonics as
\[
P(r\xi, \eta) = \frac{1 - r^2}{(1 - 2r(\xi \cdot \eta) + r^2)^{n/2}} = \sum_{k=0}^\infty r^k \frac{\omega_{n-1}}{d_k} Z_\eta^{(k)}(\xi).
\]
\end{theorem}

\begin{corollary}\label{cor:c4-2-13}
The projection operator $S_k : L^2(\Sigma) \to \mathcal{Y}_k$ is given by
\[
(S_k f)(\xi) = \int_\Sigma Z_\xi^{(k)}(\eta) f(\eta)\,d\eta.
\]
\end{corollary}

\begin{theorem}[Funk--Hecke Theorem]\label{thm:c4-2-14}
Let $K(t)$ be a function such that $\int_{-1}^1 |K(t)| (1 - t^2)^{(n-3)/2}\,dt < \infty$. If $Y_k \in \mathcal{Y}_k$, then for all $\xi \in \Sigma$,
\[
\int_\Sigma K(\xi \cdot \eta) Y_k(\eta)\,d\eta = \lambda_k Y_k(\xi),
\]
where the eigenvalue $\lambda_k$ is given by
\[
\lambda_k = \frac{\omega_{n-2}}{C_k^\lambda(1)} \int_{-1}^1 K(t) C_k^\lambda(t) (1 - t^2)^{(n-3)/2}\,dt.
\]
\end{theorem}

The expansion of the Poisson kernel stated in Theorem~\ref{thm:c4-2-12} can, therefore, be written in the equivalent form
\begin{equation}\label{eq:c4-2-15}
\frac{1 - r^2}{(1 - 2r(\xi \cdot \eta) + r^2)^{n/2}} = \sum_{k=0}^\infty r^k \frac{\omega_{n-1}}{d_k} Z_\eta^{(k)}(\xi).
\end{equation}

\begin{corollary}\label{cor:c4-2-16}
Every operator on $L^2(\Sigma)$ defined by a rotation-invariant kernel is diagonalized by the spherical harmonics.
\end{corollary}

\begin{corollary}\label{cor:c4-2-17}
The Laplace--Beltrami operator on $\Sigma$ has eigenvalues $-k(k + n - 2)$ with eigenspaces $\mathcal{Y}_k$.
\end{corollary}

\begin{lemma}\label{lem:c4-2-18}
Let $P \in \mathcal{H}_k$. Then in spherical coordinates $x = r\xi$, $\Delta(f(r) P(x)) = (\Delta_r f) P(x) + r^{-2} f(r) \Delta_\Sigma P(\xi)$.
\end{lemma}

\begin{equation}\label{eq:c4-2-19}
\Delta(f(r) Y_k(\xi)) = \left( \frac{d^2 f}{dr^2} + \frac{n-1}{r} \frac{df}{dr} - \frac{k(k + n - 2)}{r^2} f(r) \right) Y_k(\xi).
\end{equation}

\section{The Action of the Fourier Transform on the Spaces \texorpdfstring{$\mathfrak{H}_k$}{Hk}}

Let $\mathfrak{H}_k$ denote the subspace of $L^2(E_n)$ consisting of functions of the form $f(x) = f_0(|x|) P_k(x)$ with $P_k \in \mathcal{H}_k$. These spaces are invariant under rotations and mutually orthogonal.

\begin{lemma}\label{lem:c4-3-1}
If $k$ is a non-negative integer and $P_k \in \mathcal{H}_k$, then
\begin{equation}\label{eq:c4-3-2}
\int_\Sigma e^{-2\pi i r \rho (\xi \cdot \eta)} P_k(\eta)\,d\eta = 2\pi i^{-k} (r\rho)^{-\frac{n-2}{2}} J_{k + \frac{n-2}{2}}(2\pi r \rho) P_k(\xi).
\end{equation}
\end{lemma}

\begin{theorem}\label{thm:c4-3-3}
The subspace $\mathfrak{H}_k$ is invariant under the Fourier transform. More precisely, if $f(x) = f_0(|x|) P_k(x) \in \mathfrak{H}_k$, then its Fourier transform is given by $\hat{f}(y) = F_0(|y|) P_k(y)$ for a suitable radial factor $F_0$.
\end{theorem}

\begin{theorem}\label{thm:c4-3-4}
When $k = 0$, $P_0(x) \equiv 1$ and $\mathfrak{H}_0$ is the space of radial functions in $L^2(E_n)$; the Fourier transform of a radial function is radial.
\end{theorem}

The factor $F_0$ of Theorem~\ref{thm:c4-3-3} is given explicitly by
\begin{equation}\label{eq:c4-3-5}
F_0(\rho) = 2\pi i^{-k} \rho^{-\left(k + \frac{n-2}{2}\right)} \int_0^\infty f_0(r) r^{k + \frac{n}{2}} J_{k + \frac{n-2}{2}}(2\pi r \rho)\,dr,
\end{equation}
while the special case $k = 0$ of Theorem~\ref{thm:c4-3-4} takes the form
\begin{equation}\label{eq:c4-3-6}
\hat{f}(\rho) = 2\pi \rho^{-\frac{n-2}{2}} \int_0^\infty f_0(r) r^{n/2} J_{\frac{n-2}{2}}(2\pi r \rho)\,dr.
\end{equation}
A related identity, useful in the applications, is
\begin{equation}\label{eq:c4-3-7}
\int_\Sigma Z_\eta^{(k)}(\xi) P_k(\eta)\,d\eta = \frac{\omega_{n-1}}{C_k^\lambda(1)} P_k(\xi).
\end{equation}
Passing to the differential operators associated with $\mathcal{H}_k$ we have
\begin{equation}\label{eq:c4-3-8}
\widehat{P_k(D) f}(y) = P_k(2\pi i y) \hat{f}(y),
\end{equation}
and, for the Gaussian,
\begin{equation}\label{eq:c4-3-9}
P_k(D)(e^{-\pi |x|^2}) = (-2\pi)^k P_k(x) e^{-\pi |x|^2}.
\end{equation}

\begin{theorem}\label{thm:c4-3-10}
The functions $P_k(x) e^{-\pi |x|^2}$, where $P_k \in \mathcal{H}_k$, are eigenfunctions of the Fourier transform with eigenvalue $(-i)^k$:
\[
\mathcal{F}(P_k(x) e^{-\pi |x|^2}) = (-i)^k P_k(y) e^{-\pi |y|^2}.
\]
\end{theorem}

\begin{lemma}\label{lem:c4-3-11}
For $\alpha > 0$ and $\text{Re}\,\beta > -k - n$, the Fourier transform of $|x|^\beta P_k(x)$ in the sense of distributions is
\begin{equation}\label{eq:c4-3-12}
\mathcal{F}(|x|^\beta P_k(x)) = i^{-k} \pi^{-\beta - k - \frac{n}{2}} \frac{\Gamma\left(\frac{k + n + \beta}{2}\right)}{\Gamma\left(\frac{k - \beta}{2}\right)} |y|^{-\beta - k - n} P_k(y).
\end{equation}
\end{lemma}

\section{Some Applications}

\begin{theorem}[Riesz Potentials]\label{thm:c4-4-1}
For $0 < \alpha < n$, the Riesz potential $I_\alpha f$ of a function $f$ is defined by
\begin{equation}\label{eq:c4-4-2}
(I_\alpha f)(x) = \frac{1}{\gamma(\alpha)} \int_{E_n} \frac{f(t)}{|x - t|^{n - \alpha}}\,dt,
\end{equation}
where $\gamma(\alpha) = \pi^{n/2} 2^\alpha \Gamma(\alpha/2) / \Gamma((n - \alpha)/2)$. The Fourier transform of the kernel $|x|^{\alpha - n} / \gamma(\alpha)$ is $(2\pi |y|)^{-\alpha}$.
\end{theorem}

\begin{equation}\label{eq:c4-4-3}
I_\alpha(I_\beta f) = I_{\alpha + \beta} f \quad (\alpha, \beta > 0, \alpha + \beta < n).
\end{equation}

\begin{equation}\label{eq:c4-4-4}
\Delta(I_2 f) = -f.
\end{equation}

\begin{theorem}[Calder\'on--Zygmund Kernels]\label{thm:c4-4-5}
Let $P_k \in \mathcal{H}_k$ with $k \geq 1$. The kernel $K(x) = P_k(x) / |x|^{n + k}$ is a singular integral kernel of Calder\'on--Zygmund type. Its principal value Fourier transform is
\begin{equation}\label{eq:c4-4-6}
\widehat{\text{p.v.}\,K}(y) = (-i)^k \pi^{n/2} \frac{\Gamma(k/2)}{\Gamma((n + k)/2)} \frac{P_k(y)}{|y|^k}.
\end{equation}
\end{theorem}

\begin{theorem}[Riesz Transforms]\label{thm:c4-4-7}
When $k = 1$, $P_1(x) = x_j$. The singular integral operators
\begin{equation}\label{eq:c4-4-8}
(R_j f)(x) = \lim_{\varepsilon \to 0} \frac{\Gamma\left(\frac{n+1}{2}\right)}{\pi^{(n+1)/2}} \int_{|t| > \varepsilon} \frac{t_j}{|t|^{n+1}} f(x - t)\,dt
\end{equation}
are called the \emph{Riesz transforms} on $E_n$. Their Fourier multiplier is
\begin{equation}\label{eq:c4-4-9}
\widehat{R_j f}(y) = -i \frac{y_j}{|y|} \hat{f}(y), \quad j = 1, 2, \dots, n.
\end{equation}
\end{theorem}

\begin{equation}\label{eq:c4-4-10}
\sum_{j=1}^n R_j^2 = -I.
\end{equation}

\begin{theorem}\label{thm:c4-4-11}
The Riesz transforms $R_j$ are bounded linear operators on $L^p(E_n)$ for all $1 < p < \infty$, and are of weak type $(1, 1)$.
\end{theorem}

\begin{corollary}\label{cor:c4-4-12}
If $\Delta u = f$ in $E_n$, then the second partial derivatives satisfy $\partial^2 u / \partial x_j \partial x_k = -R_j R_k f$. Consequently, for $1 < p < \infty$,
\[
\left\| \frac{\partial^2 u}{\partial x_j \partial x_k} \right\|_p \leq A_p \|\Delta u\|_p.
\]
\end{corollary}

\begin{lemma}\label{lem:c4-4-13}
The Bochner--Riesz kernel of index $\delta$ is given by
\begin{equation}\label{eq:c4-4-14}
\Phi_\delta(x) = \pi^{-\delta} \Gamma(\delta + 1) |x|^{-\left(\frac{n}{2} + \delta\right)} J_{\frac{n}{2} + \delta}(2\pi |x|).
\end{equation}
\end{lemma}

\begin{theorem}\label{thm:c4-4-15}
If $\delta > (n - 1)/2$, the Bochner--Riesz means $S_R^\delta f = f * (\Phi_\delta)_R$ converge to $f$ in $L^p(E_n)$ for all $1 \leq p \leq \infty$.
\end{theorem}

\section{Further Results}

\begin{enumerate}
\item[5.1.] \textbf{Addition theorem for spherical harmonics.} The relation between Gegenbauer polynomials and zonal harmonics yields the classical addition theorem:
\[
C_k^\lambda(\xi \cdot \eta) = \frac{d_k C_k^\lambda(1)}{\omega_{n-1}} \sum_{j=1}^{d_k} Y_k^{(j)}(\xi) \overline{Y_k^{(j)}(\eta)}.
\]

\item[5.2.] \textbf{Expansion of plane waves.} The plane wave $e^{i x \cdot y}$ can be expanded in spherical harmonics and Bessel functions (the Rayleigh expansion):
\[
e^{i x \cdot y} = (2\pi)^{n/2} \sum_{k=0}^\infty i^k (|x||y|)^{-\frac{n-2}{2}} J_{k + \frac{n-2}{2}}(|x||y|) Z_{\eta}^{(k)}(\xi).
\]

\item[5.3.] \textbf{Hermite expansions.} The eigenfunctions $P_k(x) e^{-\pi |x|^2}$ can be organized into the classical multidimensional Hermite functions, forming a complete orthonormal basis of $L^2(E_n)$ consisting of eigenfunctions of the Fourier transform.

\item[5.4.] \textbf{Higher-order Riesz transforms.} The operators with multipliers $P_k(y)/|y|^k$ for $P_k \in \mathcal{H}_k$ form a commutative algebra of singular integral operators generated by the Riesz transforms $R_j$.
\end{enumerate}

\bibnoteshead

We introduced only those properties of Bessel functions that were needed for our study of the Fourier transform. For a thorough treatment of this subject see Watson~\cite{Watson1} or the Bateman Manuscript Project~\cite{Erdelyi1}, Vol.~2. We also refer the reader to the latter for further details concerning spherical harmonics, Gegenbauer and Tchebichef polynomials. For earlier studies of the action of the Fourier transform on the spaces $\mathfrak{H}_k$ see Hecke~\cite{Hecke1}, Bochner~\cite{Bochner6}, and Herz~\cite{Herz1}. For another proof of Theorem~4.1 see Calder\'on~\cite{Calderon3}. The Fourier transforms of the kernels $K(x) = \Omega(x')/|x|^n$ were first thoroughly studied by Mihlin~\cite{Mihlin1}, and Calder\'on and Zygmund~\cite{CalderonZygmund5}. The properties of the rotation group can be exploited further to obtain various results related to those in this chapter. See Vilenkin~\cite{Vilenkin1}, Coifman and G.~Weiss~\cite{CoifmanWeiss1}.
